\chapter{QUBO Formulation and QAOA Method}

Figure~\ref{fig:pipeline} summarises the method. Steps 1--3 and 5--6 run on a classical computer; step 4 is the quantum step.

\begin{figure}[H]
\centering
\small
\setlength{\fboxsep}{6pt}
\begin{tabular}{@{}c@{\,$\rightarrow$\,}c@{\,$\rightarrow$\,}c@{\,$\rightarrow$\,}c@{\,$\rightarrow$\,}c@{\,$\rightarrow$\,}c@{}}
\fbox{\parbox{1.5cm}{\centering 1.\ System data}} &
\fbox{\parbox{1.5cm}{\centering 2.\ Binary encoding}} &
\fbox{\parbox{1.5cm}{\centering 3.\ QUBO $\to$ Ising}} &
\fbox{\parbox{1.5cm}{\centering\bfseries 4.\ CVaR-QAOA}} &
\fbox{\parbox{1.5cm}{\centering 5.\ Decode \& refine}} &
\fbox{\parbox{1.5cm}{\centering 6.\ Validate}}
\end{tabular}
\caption{Hybrid quantum--classical workflow.}
\label{fig:pipeline}
\end{figure}

\section{Binary encoding of generator outputs}
Each generator output is represented by $n = 4$ binary variables:
\begin{equation}
P_i(x) = P_i^{\min} + \Delta P_i \sum_{k=0}^{n-1} 2^k x_{ik}, \qquad \Delta P_i = \frac{P_i^{\max} - P_i^{\min}}{2^n - 1}.
\label{eq:enc}
\end{equation}
With $n = 4$, each unit takes one of 16 levels from $P_i^{\min}$ to $P_i^{\max}$; every bitstring therefore satisfies the generator limits \eqref{eq:limits} automatically, and the limits need no penalty. The encoding steps are $\Delta P = 10$, $7.5$ and $9$~MW for G1, G2 and G3. Three units give 12 binary variables, mapped to 12 qubits (variable $j = 4i + k$ on qubit $j$), and $2^{12} = 4096$ candidate dispatches.

\section{QUBO objective}
The power-balance equality is added as a quadratic penalty \cite{lucas2014}:
\begin{equation}
E(x) = \sum_{i=1}^{3} \left( a_iP_i(x)^2 + b_iP_i(x) + c_i \right) + A \left( \sum_{i=1}^{3} P_i(x) - \PD \right)^2.
\label{eq:qubo}
\end{equation}
Because $P_i(x)$ is linear in $x$, both terms are quadratic in the binary variables. Using $x_j^2 = x_j$, \eqref{eq:qubo} is written as
\begin{equation}
E(x) = c_0 + \sum_j q_j x_j + \sum_{j<l} U_{jl} x_j x_l,
\end{equation}
where the constant $c_0$, the linear weights $q_j$ and the couplings $U_{jl}$ are computed from the generator data and $\PD$. Each period has its own QUBO, because $\PD$ changes.

\subsection{Choice of penalty weight}
If $A$ is too small, the cheapest bitstring under-supplies the demand; if $A$ is too large, the balance term dominates and the cost information is lost in a landscape of nearly equal energies. The optimum of the QUBO was found by exhaustive search for $A \in \{1, 5, 20, 100\}$ in all 12 periods (Table~\ref{tab:penalty}).

\begin{table}[H]
\centering
\caption{Penalty sweep: QUBO optimum over all 12 periods.}
\label{tab:penalty}
\begin{tabular}{crr}
\toprule
$A$ (\$/MW$^2$h) & Max.\ $|\sum P - \PD|$ (MW) & Mean cost difference vs SLSQP\\
\midrule
1 & 7.0 & $-2.00$\,\%\\
5 & 2.0 & $-0.39$\,\%\\
\textbf{20} & \textbf{0.5} & $\mathbf{+0.011}$\,\%\\
100 & 0.0 & $+0.176$\,\%\\
\bottomrule
\end{tabular}
\end{table}

A negative cost difference is not a saving: it comes from supplying less than the demand. $A = 20$ is the smallest value that keeps the mismatch at or below 0.5~MW in every period and is used throughout.

\section{Ising Hamiltonian}
The substitution $x_j = (1 - Z_j)/2$ turns the QUBO into the Ising Hamiltonian
\begin{equation}
H_C = \sum_j h_j Z_j + \sum_{j<l} J_{jl} Z_j Z_l, \qquad
h_j = -\frac{q_j}{2} - \frac{1}{4}\sum_{l \ne j} U_{jl}, \qquad
J_{jl} = \frac{U_{jl}}{4},
\end{equation}
with a constant offset $c_0 + \tfrac12\sum_j q_j + \tfrac14\sum_{j<l}U_{jl}$, so that $E(x) = \langle x|H_C|x\rangle + \text{offset}$ for every computational basis state $|x\rangle$. This identity was checked for all 4096 bitstrings. Since the encoding couples every pair of bits through the balance penalty, $H_C$ has 12 linear terms and 66 two-body terms. For training, $H_C$ is divided by its largest coefficient magnitude, which keeps the useful range of $\gamma$ near $[0, \pi]$.

\section{The QAOA circuit}
QAOA with depth $p$ prepares
\begin{equation}
|\psi(\boldsymbol\gamma, \boldsymbol\beta)\rangle = \prod_{\ell=1}^{p} e^{-i\beta_\ell H_M} e^{-i\gamma_\ell H_C} \, |+\rangle^{\otimes 12}, \qquad H_M = \sum_j X_j.
\end{equation}
The circuit is built from native gates:
\begin{itemize}
\item a Hadamard gate on every qubit creates $|+\rangle^{\otimes 12}$, the uniform superposition of all 4096 dispatches;
\item the cost unitary is a product of $R_Z(2\gamma_\ell h_j)$ on each qubit and $R_{ZZ}(2\gamma_\ell J_{jl})$ on each coupled pair, exact because all terms of $H_C$ commute;
\item the mixer unitary is $R_X(2\beta_\ell)$ on every qubit;
\item all qubits are measured in the computational basis.
\end{itemize}
The circuit has $2p$ parameters.

\section{Training}
\subsection{CVaR objective}
Instead of the mean energy $\langle H_C \rangle$, the parameters minimise the conditional value-at-risk \cite{barkoutsos2020}: the mean energy of the lowest-energy fraction $\alpha$ of the output distribution. With the basis states sorted by energy $E_{(1)} \le E_{(2)} \le \dots$ and probabilities $p_{(k)}$,
\begin{equation}
\mathrm{CVaR}_\alpha = \frac{1}{\alpha} \sum_k E_{(k)} \, \min\!\left( p_{(k)},\ \max\!\left(0,\ \alpha - \textstyle\sum_{m<k} p_{(m)}\right) \right).
\end{equation}
$\alpha = 0.05$ is used. Minimising CVaR rewards placing probability on the best dispatches rather than improving the average one, which suits the goal of sampling the optimum. The value is normalised to $[0, 1]$ between the minimum and maximum QUBO energy for reporting.

\subsection{Optimiser and initialisation}
The classical optimiser is COBYLA (gradient-free, at most 400 iterations). At $p = 1$, eight random starting points are used and the best result is kept. Depth $p + 1$ is initialised from the optimum at depth $p$ by the INTERP rule of Zhou et al. \cite{zhou2020},
\begin{equation}
\gamma^{(p+1)}_i = \frac{i-1}{p}\gamma^{(p)}_{i-1} + \frac{p-i+1}{p}\gamma^{(p)}_{i}, \qquad i = 1,\dots,p+1,
\end{equation}
with $\gamma^{(p)}_0 = \gamma^{(p)}_{p+1} = 0$, and the same rule for $\boldsymbol\beta$. Training uses the exact statevector of the noise-free simulator.

\section{Sampling, decoding and hybrid refinement}
The trained circuit is measured 4096 times on the Qiskit Aer simulator. Each distinct bitstring is decoded to MW with \eqref{eq:enc}, and the five lowest-energy bitstrings are kept. Because the grid is 7.5--10~MW wide, the best grid point usually differs slightly from the continuous optimum and may miss $\PD$ by up to 0.5~MW. Each of the five candidates is therefore refined: every generator may move at most one grid step $\pm\Delta P_i$ from its QAOA level, and the small continuous problem inside this box is solved with SLSQP to meet $\PD$ exactly. The cheapest refined dispatch is the hybrid result. The box is widened only if demand cannot be met inside it; this was never needed for the lossless case.

The refinement does not repeat the whole optimisation: it only searches a box of about one grid cell chosen by QAOA. Its quality therefore depends on QAOA choosing the right cell, which is measured in Chapter~6.

\section{Performance measures}
\begin{itemize}
\item \textbf{Cost gap}: $(C_{\text{method}} - C_{\text{SLSQP}})/C_{\text{SLSQP}} \times 100\,\%$.
\item \textbf{Mismatch}: $\sum_i P_i - \PD$ in MW.
\item \textbf{Amplification}: probability of measuring the QUBO-optimal bitstring divided by its probability under uniform sampling ($1/4096$).
\item \textbf{Near-optimal probability}: probability that a measured bitstring is a usable starting point, i.e.\ its dispatch is within one grid step of $\PD$ and its cost within 1\,\% of the SLSQP cost.
\item \textbf{Success}: a run succeeds if its final cost is within 0.1\,\% of SLSQP and $|\text{mismatch}| \le 0.01$~MW.
\end{itemize}
